\documentclass[11pt,preprint]{elsarticle}

\usepackage{lmodern}
%%%% My spacing
\usepackage{setspace}
\setstretch{1.2}
\DeclareMathSizes{12}{14}{10}{10}

% Wrap around which gives all figures included the [H] command, or places it "here". This can be tedious to code in Rmarkdown.
\usepackage{float}
\let\origfigure\figure
\let\endorigfigure\endfigure
\renewenvironment{figure}[1][2] {
    \expandafter\origfigure\expandafter[H]
} {
    \endorigfigure
}

\let\origtable\table
\let\endorigtable\endtable
\renewenvironment{table}[1][2] {
    \expandafter\origtable\expandafter[H]
} {
    \endorigtable
}


\usepackage{ifxetex,ifluatex}
\usepackage{fixltx2e} % provides \textsubscript
\ifnum 0\ifxetex 1\fi\ifluatex 1\fi=0 % if pdftex
  \usepackage[T1]{fontenc}
  \usepackage[utf8]{inputenc}
\else % if luatex or xelatex
  \ifxetex
    \usepackage{mathspec}
    \usepackage{xltxtra,xunicode}
  \else
    \usepackage{fontspec}
  \fi
  \defaultfontfeatures{Mapping=tex-text,Scale=MatchLowercase}
  \newcommand{\euro}{€}
\fi

\usepackage{amssymb, amsmath, amsthm, amsfonts}

\def\bibsection{\section*{References}} %%% Make "References" appear before bibliography


\usepackage[numbers]{natbib}

\usepackage{longtable}
\usepackage[margin=2.3cm,bottom=2cm,top=2.5cm, includefoot]{geometry}
\usepackage{fancyhdr}
\usepackage[bottom, hang, flushmargin]{footmisc}
\usepackage{graphicx}
\numberwithin{equation}{section}
\numberwithin{figure}{section}
\numberwithin{table}{section}
\setlength{\parindent}{0cm}
\setlength{\parskip}{1.3ex plus 0.5ex minus 0.3ex}
\usepackage{textcomp}
\renewcommand{\headrulewidth}{0.2pt}
\renewcommand{\footrulewidth}{0.3pt}

\usepackage{array}
\newcolumntype{x}[1]{>{\centering\arraybackslash\hspace{0pt}}p{#1}}

%%%%  Remove the "preprint submitted to" part. Don't worry about this either, it just looks better without it:

 \def\tightlist{} % This allows for subbullets!

\usepackage{hyperref}
\hypersetup{breaklinks=true,
            bookmarks=true,
            colorlinks=true,
            citecolor=blue,
            urlcolor=blue,
            linkcolor=blue,
            pdfborder={0 0 0}}


% The following packages allow huxtable to work:
\usepackage{siunitx}
\usepackage{multirow}
\usepackage{hhline}
\usepackage{calc}
\usepackage{tabularx}
\usepackage{booktabs}
\usepackage{caption}


\newenvironment{columns}[1][]{}{}

\newenvironment{column}[1]{\begin{minipage}{#1}\ignorespaces}{%
\end{minipage}
\ifhmode\unskip\fi
\aftergroup\useignorespacesandallpars}

\def\useignorespacesandallpars#1\ignorespaces\fi{%
#1\fi\ignorespacesandallpars}

\makeatletter
\def\ignorespacesandallpars{%
  \@ifnextchar\par
    {\expandafter\ignorespacesandallpars\@gobble}%
    {}%
}
\makeatother


% definitions for citeproc citations
\NewDocumentCommand\citeproctext{}{}
\NewDocumentCommand\citeproc{mm}{%
\href{\#cite.\detokenize{#1}}{#2}\nocite{#1}}

\makeatletter
% allow citations to break across lines
\let\@cite@ofmt\@firstofone
% avoid brackets around text for \cite:
\def\@biblabel#1{}
\def\@cite#1#2{{#1\if@tempswa , #2\fi}}
\makeatother
\newlength{\cslhangindent}
\setlength{\cslhangindent}{1.5em}
\newlength{\csllabelwidth}
\setlength{\csllabelwidth}{3em}
\newenvironment{CSLReferences}[2] % #1 hanging-indent, #2 entry-spacing
{\begin{list}{}{%
	\setlength{\itemindent}{0pt}
	\setlength{\leftmargin}{0pt}
	\setlength{\parsep}{0pt}
	% turn on hanging indent if param 1 is 1
	\ifodd #1
	\setlength{\leftmargin}{\cslhangindent}
	\setlength{\itemindent}{-1\cslhangindent}
	\fi
	% set entry spacing
	\setlength{\itemsep}{#2\baselineskip}}}
{\end{list}}

\usepackage{calc}
\newcommand{\CSLBlock}[1]{\hfill\break\parbox[t]{\linewidth}{\strut\ignorespaces#1\strut}}
\newcommand{\CSLLeftMargin}[1]{\parbox[t]{\csllabelwidth}{\strut#1\strut}}
\newcommand{\CSLRightInline}[1]{\parbox[t]{\linewidth - \csllabelwidth}{\strut#1\strut}}
\newcommand{\CSLIndent}[1]{\hspace{\cslhangindent}#1}


\urlstyle{same}  % don't use monospace font for urls
\setlength{\parindent}{0pt}
\setlength{\parskip}{6pt plus 2pt minus 1pt}
\setlength{\emergencystretch}{3em}  % prevent overfull lines
\setcounter{secnumdepth}{5}

%%% Use protect on footnotes to avoid problems with footnotes in titles
\let\rmarkdownfootnote\footnote%
\def\footnote{\protect\rmarkdownfootnote}
\IfFileExists{upquote.sty}{\usepackage{upquote}}{}

%%% Include extra packages specified by user

%%% Hard setting column skips for reports - this ensures greater consistency and control over the length settings in the document.
%% page layout
%% paragraphs
\setlength{\baselineskip}{12pt plus 0pt minus 0pt}
\setlength{\parskip}{12pt plus 0pt minus 0pt}
\setlength{\parindent}{0pt plus 0pt minus 0pt}
%% floats
\setlength{\floatsep}{12pt plus 0 pt minus 0pt}
\setlength{\textfloatsep}{20pt plus 0pt minus 0pt}
\setlength{\intextsep}{14pt plus 0pt minus 0pt}
\setlength{\dbltextfloatsep}{20pt plus 0pt minus 0pt}
\setlength{\dblfloatsep}{14pt plus 0pt minus 0pt}
%% maths
\setlength{\abovedisplayskip}{12pt plus 0pt minus 0pt}
\setlength{\belowdisplayskip}{12pt plus 0pt minus 0pt}
%% lists
\setlength{\topsep}{10pt plus 0pt minus 0pt}
\setlength{\partopsep}{3pt plus 0pt minus 0pt}
\setlength{\itemsep}{5pt plus 0pt minus 0pt}
\setlength{\labelsep}{8mm plus 0mm minus 0mm}
\setlength{\parsep}{\the\parskip}
\setlength{\listparindent}{\the\parindent}
%% verbatim
\setlength{\fboxsep}{5pt plus 0pt minus 0pt}



\begin{document}



\begin{frontmatter}  %

\title{Loans and Credit}

% Set to FALSE if wanting to remove title (for submission)




\author[Add1]{25424971}
\ead{25424971@sun.ac.za}








\vspace{1cm}





\vspace{0.5cm}

\end{frontmatter}

\setcounter{footnote}{0}



%________________________
% Header and Footers
%%%%%%%%%%%%%%%%%%%%%%%%%%%%%%%%%
\pagestyle{fancy}
\chead{}
\rhead{}
\lfoot{}
\rfoot{\footnotesize Page \thepage}
\lhead{}
%\rfoot{\footnotesize Page \thepage } % "e.g. Page 2"
\cfoot{}

%\setlength\headheight{30pt}
%%%%%%%%%%%%%%%%%%%%%%%%%%%%%%%%%
%________________________

\headsep 35pt % So that header does not go over title




\section{Introduction}\label{introduction}

The key drivers of default risk are explored. The following factors of
individuals are considered: whether the income source was verified, the
debt-to-income ratio, the number of year that the individual has had any
sort of credit line, prior bankruptcies, the number of mortgage
accounts, bankcard utilisation, interest rate, term length, and history
of credit behaviour.

\subsection{The Key Risk Drivers for
Default}\label{the-key-risk-drivers-for-default}

A logistic regression is used to study the determinants of default risk,
as the dependent variable is binary in nature. Across both
specifications, all the coefficients are statistically significant,
justifying their statuses as key drivers of default risk. The Lending
Club's credit grade system proves to be great indicator of debt default,
with lower grades being associated with an increased probability of
default. The DTI ratio, prior bankruptcies, and bankcard utilisation are
also associated with increased default risk. The number of mortgage
accounts and never having missed payments on debt are associated with
lower default risk.

\begin{verbatim}
## 
## Logistic Regression: Predictors of Loan Default
## =====================================================
##                             Dependent variable:      
##                        ------------------------------
##                         Default (1 = Loss, 0 = Paid) 
##                              (1)            (2)      
## -----------------------------------------------------
## Source Verified           0.346***        0.188***   
##                            (0.010)        (0.011)    
## Verified                  0.555***        0.283***   
##                            (0.011)        (0.012)    
## DTI                       0.027***        0.015***   
##                           (0.0005)        (0.0005)   
## Credit Age                -0.006***       0.004***   
##                            (0.001)        (0.001)    
## Prior Bankruptcies        0.179***        0.079***   
##                            (0.009)        (0.010)    
## % Never Delinquent        -0.007***      -0.004***   
##                           (0.0004)        (0.0005)   
## Mortgage Accounts         -0.126***      -0.123***   
##                            (0.003)        (0.003)    
## Bankcard Utilisation      0.005***        0.002***   
##                           (0.0001)        (0.0002)   
## No. Revolving Accounts    0.005***        0.005***   
##                            (0.001)        (0.001)    
## Grade B                                   0.655***   
##                                           (0.020)    
## Grade C                                   1.044***   
##                                           (0.026)    
## Grade D                                   1.290***   
##                                           (0.037)    
## Grade E                                   1.413***   
##                                           (0.048)    
## Grade F                                   1.590***   
##                                           (0.060)    
## Grade G                                   1.678***   
##                                           (0.078)    
## Interest Rate                             0.028***   
##                                           (0.003)    
## Term                                      0.016***   
##                                           (0.0004)   
## Constant                  -1.409***      -3.419***   
##                            (0.047)        (0.056)    
## -----------------------------------------------------
## Observations               372,684        372,684    
## Log Likelihood          -191,565.700    -180,892.700 
## Akaike Inf. Crit.        383,151.300    361,821.500  
## =====================================================
## Note:                   *p<0.05; **p<0.01; ***p<0.001
\end{verbatim}

\subsection{Default-Risk Customer
Profiles}\label{default-risk-customer-profiles}

\subsubsection{Default Rate by Lending Club's Customer Grading
System}\label{default-rate-by-lending-clubs-customer-grading-system}

The plot below illustrates that default rates increase as credit grades
worsen. Customers that are assigned poor Lending Club credit grades are
at higher risk of defaulting on their debt.

\begin{figure}[H]

{\centering \includegraphics{25424971Question3_files/figure-latex/unnamed-chunk-4-1} 

}

\caption{Default Rate by Grading System .\label{Figure1}}\label{fig:unnamed-chunk-4}
\end{figure}

\subsubsection{Default Rate by Home
Ownership}\label{default-rate-by-home-ownership}

Individuals that rent instead of own their house/mortgage have an
average default rate of 26.9\%. Owning a mortgage is associated with the
lowest average default rate, supporting the findings of the logit
regression, where the number of mortgages negatively relates to the
default rate. This may simply be a proxy for income.

\begin{figure}[H]

{\centering \includegraphics{25424971Question3_files/figure-latex/unnamed-chunk-5-1} 

}

\caption{Default Rate by Home Ownership.\label{Figure1}}\label{fig:unnamed-chunk-5}
\end{figure}

\subsubsection{Default Rate by Employment
Length}\label{default-rate-by-employment-length}

The default rate appers to be similarly distributed across employment
lengths, with a slight downward trend, such that longer employment years
(stable income) indicate a slightly lower risk of default.

\begin{figure}[H]

{\centering \includegraphics{25424971Question3_files/figure-latex/unnamed-chunk-6-1} 

}

\caption{Default Rate by Employment Length.\label{Figure1}}\label{fig:unnamed-chunk-6}
\end{figure}

\subsubsection{Default Rate by State for Both Long and Short Term
Loans}\label{default-rate-by-state-for-both-long-and-short-term-loans}

Arkansas, Los Angeles, Mississipi, Nebraska, and Oklahoma have the
highest average default rates, signalling that individuals from these
states iintroduce increased risk. Whereas individuals from DC, Maine,
New Hampshire, Oregan, Vermont, and West Virginia signal prudence, as
these states have the lowest average default rates.

\begin{figure}[H]

{\centering \includegraphics{25424971Question3_files/figure-latex/unnamed-chunk-7-1} 

}

\caption{Default Rate by State for Both Long and Short Term Loans.\label{Figure1}}\label{fig:unnamed-chunk-7}
\end{figure}

\subsubsection{Default Rate by State for Short Term Loans
Only}\label{default-rate-by-state-for-short-term-loans-only}

The state analysis is similar when only short-term loans are considered.

\begin{figure}[H]

{\centering \includegraphics{25424971Question3_files/figure-latex/unnamed-chunk-8-1} 

}

\caption{Default Rate by State for Short Term Loans Only.\label{Figure1}}\label{fig:unnamed-chunk-8}
\end{figure}

\subsubsection{Default Rate by Purpose of
Loan}\label{default-rate-by-purpose-of-loan}

Individuals who cite funding a small business or planning a wedding as
the reason for their debt obligations are at the highest risk of
defaulting on their loans.

\begin{figure}[H]

{\centering \includegraphics{25424971Question3_files/figure-latex/unnamed-chunk-9-1} 

}

\caption{Default Rate by Purpose.\label{Figure1}}\label{fig:unnamed-chunk-9}
\end{figure}

\subsection{The Debt-to-Income and Loan Status
Relationship}\label{the-debt-to-income-and-loan-status-relationship}

The average debt-to-income ratio is the highest for those who have
defaulted on their loans and completely written their loans off,
although the variance is not large across these broader loan status
groups.

\begin{figure}[H]

{\centering \includegraphics{25424971Question3_files/figure-latex/unnamed-chunk-10-1} 

}

\caption{Average DTI by Loan Status.\label{Figure1}}\label{fig:unnamed-chunk-10}
\end{figure}

\section{Conclusion}\label{conclusion}

Multiple factors must be considered when approving loan applications to
ensure that the default rate is minimised. By scrutinizing over where
individuals reside, the debt-to-income ratios, the purpose of the
credit, the credit category, home ownership, length of employment, debt
history, and bankcard utilisation, debt default risk can be improved.

\bibliography{Tex/ref}





\end{document}
